\chapter{اجرا و بررسی نتیجه}

\begin{goalbox}
ربات را روشن کنید، یک گفت‌وگوی واقعی در تلگرام انجام دهید و آزمون‌های آفلاین را اجرا کنید.
\end{goalbox}

مطمئن شوید محیط مجازی فعال و متغیر
\lr{BOT\_TOKEN}
تنظیم شده است؛ سپس اجرا کنید:

\begin{LTR}
\begin{lstlisting}[language=bash,numbers=none]
python bot.py
\end{lstlisting}

\begin{resultbox}
\texttt{Bot is running. Press Ctrl+C to stop it.}
\end{resultbox}
\end{LTR}

ترمینال را باز نگه دارید. در تلگرام پیوند داده‌شده توسط بات‌فادر را باز کنید یا
نام کاربری ساخته‌شده را جست‌وجو کنید. دکمهٔ شروع را بزنید.

\section{بررسی دستی در چهار بخش}

\begin{enumerate}
\item دستور
\lr{/start}
را بفرستید. پاسخ با نام کوچک شما در تلگرام آغاز می‌شود و هر دو دستور را نشان می‌دهد.
\item دستور
\lr{/help}
را بفرستید. پاسخ سه گام کوتاه استفاده را نمایش می‌دهد.
\item پیام
\lr{My first bot}
را بفرستید. پاسخ باید دقیقاً عبارت زیر باشد:
\lr{You said: My first bot}
\item یک عکس بفرستید. پاسخ باید دقیقاً عبارت زیر باشد:
\lr{Please send a text message so I can repeat it.}
\end{enumerate}

اگر هر چهار نتیجه یکسان بود، مسیر کامل شامل تلگرام، توکن، پایتون، تله‌بات و
پاسخ‌دهنده‌ها درست کار می‌کند.

\section{اجرای آزمون‌های آفلاین}

ربات را با
\lr{Ctrl+C}
متوقف کنید؛ سپس اجرا کنید:

\begin{LTR}
\begin{lstlisting}[language=bash,numbers=none]
python -m unittest -v
\end{lstlisting}
\end{LTR}

\begin{resultbox}
چهار آزمون فهرست می‌شوند و سپس عبارت
\lr{OK}
نمایش داده می‌شود. این بررسی‌ها به توکن نیاز ندارند و به تلگرام متصل نمی‌شوند؛
آن‌ها تابع‌های سازندهٔ پاسخ را بررسی می‌کنند.
\end{resultbox}

هر زمان پیامی را در فایل
\lr{bot.py}
تغییر دادید، آزمون متناظر را نیز به‌روز و این دستور را دوباره اجرا کنید.

\begin{checkpointbox}
ربات شما به هر چهار بررسی دستی پاسخ درست می‌دهد و چهار آزمون آفلاین موفق هستند.
\end{checkpointbox}
